% Herramientas de layout para el informe Quarto -> PDF.
% La idea es concentrar aqui los ajustes finos y evitar repetirlos en el .qmd.

\usepackage{array}
\usepackage{booktabs}
\usepackage{makecell}
\usepackage{tabularx}
\usepackage{adjustbox}
\usepackage{threeparttable}
\usepackage{pdflscape}
\usepackage{multicol}
\usepackage{enumitem}
\usepackage{placeins}
\usepackage{microtype}

\raggedbottom
\setlength{\emergencystretch}{3em}
\setlength{\columnsep}{0.72cm}
\setlist{nosep,leftmargin=*}

% Espaciado vertical puntual.
\newcommand{\compactspace}{\vspace{-0.6em}}
\newcommand{\verycompactspace}{\vspace{-1.0em}}
\newcommand{\openspace}{\vspace{0.6em}}
\newcommand{\keepnext}[1][6]{\Needspace{#1\baselineskip}}
\newcommand{\stopfloats}{\FloatBarrier}

% Bloque de texto mas compacto, util para paginas densas o anexos.
\newenvironment{compacttext}
  {\begingroup\small\setlength{\parskip}{0.25em}\setlength{\parindent}{0pt}}
  {\par\endgroup}

% Doble columna local. A diferencia de \twocolumn, no cambia todo el documento.
\newenvironment{twocolblock}
  {\begin{multicols}{2}\raggedcolumns}
  {\end{multicols}}

% Pagina horizontal para tablas o figuras anchas.
\newenvironment{landscapeblock}
  {\begin{landscape}\begingroup}
  {\par\endgroup\end{landscape}}

% Ajustes consistentes para tablas densas.
\newenvironment{compacttable}[1][\small]
  {\begingroup#1\setlength{\tabcolsep}{3.5pt}\renewcommand{\arraystretch}{0.86}}
  {\par\endgroup}

% Caja para escalar tablas tabular/tabularx al ancho disponible.
% No usar con longtable, porque longtable necesita partir entre paginas.
\newenvironment{fittable}[1][\textwidth]
  {\begin{adjustbox}{max width=#1}}
  {\end{adjustbox}}

% Columnas practicas para tablas manuales.
\newcolumntype{L}[1]{>{\raggedright\arraybackslash}p{#1}}
\newcolumntype{C}[1]{>{\centering\arraybackslash}p{#1}}
\newcolumntype{R}[1]{>{\raggedleft\arraybackslash}p{#1}}

% Atajos para captions mas pegados cuando una figura/tabla queda muy separada.
\newcommand{\tightcaptiongap}{\vspace{-0.55em}}
\newcommand{\widefig}[2]{%
  \begin{figure}[H]
  \centering
  \includegraphics[width=#1\textwidth]{#2}
  \end{figure}
}
